\begingroup\small\setstretch{1.04}
\section*{Sources and further reading}
\addcontentsline{toc}{section}{Sources and further reading}
\subsection*{Project sources}
\begin{enumerate}[label={[P\arabic*]},leftmargin=2.5em]
\item Komorebi P. \emph{Projector Refinement Physics, Volume I}. The source studies cover weak fields, holonomy, Gowdy reduction, linearized ADM constraints, nonlinear transfer, resolved connections, the Lorentzian parent, regular reduction, electrovac geometry, and scalar collapse.
\item Komorebi P. \emph{Projector Refinement Physics, Volume I: Technical Supplement}. The nonlinear and strong-field claim dossiers specify boundary, spectral, and numerical hypotheses. The verification appendix maps these hypotheses to the constructions used here.
\item Komorebi P. \emph{Komorebi Physics formal source}. Its projector, gauge, and free Dirac results provide related mathematical infrastructure. They do not constitute a compiler certificate for the nonlinear gravity and collapse statements.
\end{enumerate}

\subsection*{Primary context}
\begin{enumerate}[label={[C\arabic*]},leftmargin=2.5em]
\item R. Arnowitt, S. Deser, and C. W. Misner. ``The Dynamics of General Relativity.'' In L. Witten, ed., \emph{Gravitation: An Introduction to Current Research}, Wiley, 1962, Chapter 7, pp. 227--265. Author republication arXiv:gr-qc/0405109v1, 19 May 2004. \url{https://arxiv.org/abs/gr-qc/0405109v1}.
\item A. Ashtekar. ``New Variables for Classical and Quantum Gravity.'' \emph{Physical Review Letters} \textbf{57} (1986), 2244--2247. \url{https://doi.org/10.1103/PhysRevLett.57.2244}.
\item I. A. Batalin and G. A. Vilkovisky. ``Gauge algebra and quantization.'' \emph{Physics Letters B} \textbf{102} (1981), 27--31. \url{https://doi.org/10.1016/0370-2693(81)90205-7}.
\item A. S. Cattaneo, P. Mnev, and N. Reshetikhin. ``Classical BV theories on manifolds with boundary.'' \emph{Communications in Mathematical Physics} \textbf{332} (2014), 535--603. Author preprint arXiv:1201.0290v3, 28 June 2012. \url{https://arxiv.org/abs/1201.0290v3}; \url{https://doi.org/10.1007/s00220-014-2145-3}.
\item A. S. Cattaneo, P. Mnev, and N. Reshetikhin. ``Perturbative quantum gauge theories on manifolds with boundary.'' \emph{Communications in Mathematical Physics} \textbf{357} (2018), 631--730. Author preprint arXiv:1507.01221v2, 30 May 2016. \url{https://arxiv.org/abs/1507.01221v2}; \url{https://doi.org/10.1007/s00220-017-3031-6}.
\item R. I. McLachlan, G. R. W. Quispel, and N. Robidoux. ``Geometric integration using discrete gradients.'' \emph{Philosophical Transactions of the Royal Society A} \textbf{357} (1999), 1021--1045. \url{https://doi.org/10.1098/rsta.1999.0363}.
\item C. W. Misner and D. H. Sharp. ``Relativistic Equations for Adiabatic, Spherically Symmetric Gravitational Collapse.'' \emph{Physical Review} \textbf{136} (1964), B571--B576. \url{https://doi.org/10.1103/PhysRev.136.B571}.
\item J. Ziprick and G. Kunstatter. ``Numerical study of black-hole formation in Painlev\'e--Gullstrand coordinates.'' \emph{Physical Review D} \textbf{79} (2009), 101503(R). Author preprint \emph{Spherically Symmetric Black Hole Formation in Painlev\'e--Gullstrand Coordinates}, arXiv:0812.0993v1, 4 December 2008. \url{https://arxiv.org/abs/0812.0993v1}; \url{https://doi.org/10.1103/PhysRevD.79.101503}.
\end{enumerate}

\subsection*{How these sources are used}
The ADM republication supplies the canonical setting for lapse, shift, and constraints. Ashtekar and Batalin--Vilkovisky supply the historical connection and master-equation settings. These references do not establish the projector model's continuum convergence or its quantum anomaly cancellation.

The two Cattaneo--Mnev--Reshetikhin preprints provide the precise comparison for boundary reduction and BV elimination. Section 3.1.3 of arXiv:1201.0290v3 constructs boundary data through a presymplectic quotient and its projectability conditions. Bulk gauge invariance alone therefore does not discharge the missing regular boundary construction in the present gravity model. In arXiv:1507.01221v2, Theorem 2.9 and Corollary 2.12 give BV pushforward and preservation of the quantum master equation for compatible finite-dimensional odd-symplectic data. The identity is conditional on the fine quantum master equation and well-defined integration. It does not manufacture an anomaly-free quantum gravity measure.

McLachlan--Quispel--Robidoux is the established discrete-gradient context. The preservation identity used here is proved directly from the segment integral and skewness of the endpoint Dirac tensor. Misner--Sharp identifies the spherical mass tradition. The charged mass-coordinate identity in this volume is independently differentiated in the stated canonical convention. Ziprick--Kunstatter is prior work on scalar collapse in Painlev\'e--Gullstrand coordinates. Its regular-center and critical-collapse analyses are distinct from the excised, first-event benchmark retained here.

Relevant full-text sections of the ADM, both Cattaneo--Mnev--Reshetikhin, and Ziprick--Kunstatter author preprints were inspected in their stated versions. For Ashtekar, Batalin--Vilkovisky, McLachlan--Quispel--Robidoux, and Misner--Sharp, bibliographic metadata was authenticated, but publisher full-text access was unavailable. Their uninspected results are not used to discharge a proof obligation. The source record describes this focused comparison and its limits. No priority claim or exhaustive novelty review is made.

\endgroup
